\section{Introducció}

\subsection{Justificació del tema}
S’ha escollit el món de les motocicletes, ja que sempre m’han agradat molt i m’ha interessat molt la mecànica de la qual disposen. A més, he escollit la marca japonesa Honda pel fet que em sembla una marca molt interessant i amb moltes possibilitats d’aprendre com dissenyar un bon motor. Una part de la meva motivació personal és que disposo d’una motocicleta Honda a casa i tinc molt d'interès a aprendre el seu funcionament intern. Finalment, l’objectiu és trobar una motocicleta ideal per a ús personal.

\subsection{Plantejament del problema}
El plantejament d’aquesta pregunta es fa amb l’objectiu de comprovar quina és la millor motocicleta nua del segment mitjà d’Honda mitjançant criteris tècnics objectius i ponderats per poder escollir una opció òptima per a ús personal.

\subsection{Hipòtesi}
La CB650R presenta el millor rendiment global dins del segment de motocicletes nues (*naked*) de gamma mitjana quan s'avalua mitjançant criteris tècnics objectius i ponderats?

\subsection{Objectius}
L’objectiu general del treball de recerca és determinar si la CB650R és la millor motocicleta nua del segment mitjà d’Honda, comparant-la amb altres models del mateix segment. Tanmateix, es pretén augmentar els coneixements mecànics sobre les motocicletes per facilitar l'elecció d'un vehicle per a ús personal.

\subsection{Metodologia}
Per aconseguir uns resultats objectius, es realitzaran enquestes a usuaris amb la finalitat d’obtenir dades quantitatives. Seguidament, es durà a terme una entrevista a usuaris de la CB650R i, per finalitzar, es realitzarà una mitjana ponderada amb les característiques principals de les motocicletes (consum, potència i preu, entre d’altres).

\subsection{Definició del segment}
L’estudi se centrarà en motocicletes de la categoria nua (*naked*), de gamma mitjana d'Honda, d'entre 500 i 800 centímetres cúbics. S'han descartat altres segments com el turisme esportiu, l'aventura o el turisme per tenir objectius d'ús diferents. El rang de cubicatge s'estableix per garantir una comparativa justa, evitant la disparitat tècnica entre motors de cilindrades molt allunyades.

\subsection{Models a comparar}
S'han seleccionat tres models que compleixen els requisits establerts i representen els èxits de vendes d'Honda pel seu equilibri entre potència, confort i preu:

\begin{itemize}
    \item \textbf{Honda CB650R:} Motocicleta tetracilíndrica de 650 $cm^3$ i 95 CV. Presenta una estètica que fusiona l'estil clàssic (*retro*) amb l'actual.
    \item \textbf{Honda CB750 Hornet:} Motocicleta bicilíndrica de 750 $cm^3$ i 92 CV, amb un disseny aerodinàmic agressiu.
    \item \textbf{Honda CB500 Hornet:} Model bicilíndric de 500 $cm^3$ i 47 CV, optimitzat per a la conducció urbana i l'accés amb el carnet A2.
\end{itemize}
